% Cote 119: offprints, reports, courses, CV and letters (typed and printed part): archive cote 119, site page 223 (PDF page 224).
% typescript. Draft compared: none.
% AI-transcribed by Claude Opus 5.5 (claude-opus-5-5, adaptive thinking, effort high), Message Batches API, batch r119_full (msgbatch_01DwZNm5Y7YsKBWyAN4sFs8w), 2026-11-13; not yet checked by eye.
\scanpage{119/223}
\mspage{. 10 .}

vectoriels sur des topos annelés quelconques. Dans
[13 bis], j'établis, à peu de choses près, que pour un
groupe discret $G$ , la théorie des représentations
linéaires de $G$ (sur un anneau de base quelconque)
ne dépend que du complété profini'' $\hat{G}$ de $G$ .

8. \uline{Groupes formels} ([17] [16 SGA 7] [14 bis]).

C'est un sujet qui relève à la fois de la
théorie des groupes, de celle des groupes de Lie, de
la géométrie algébrique, de l'arithmétique, et (sous
la forme voisine des groupes de Barsotti-Tate) de la
théorie des systèmes locaux. Ici encore, la théorie
des schémas permet une grande aisance, et c'est dans
ce contexte par exemple que se place d'emblée I. Manin
[$*$a] , dans son exposé classique de la théorie de
Dieudonné. Ma principale contribution, en dehors de
cette simplification conceptuelle, a été le développement
d'une ``théorie de Dieudonné'' pour les groupes de
Barsotti-Tate sur des schémas de base généraux à
caractéristiques résiduelles $> 0$ , en termes du
``cristal de Dieudonné'' associé à un tel groupe. Une
esquisse de cette théorie a été exposée dans divers
cours et séminaires, y compris dans mon cours au
Collège de France en 1970/71 et 71/72 ; certains énoncés
principaux sont esquissés dans les C.R. du Congrès
International de Nice en 1970 [14 bis]. Une partie
de ces idées est développée dans la thèse de W. Messing
[$*$], et les besoins techniques de la théorie ont été
la motivation pour le développement par L. Illusie [$*$]
de sa théorie des déformations des schémas en groupes
commutatifs, vérifiant des conjectures suggérées par
cette ``théorie de Dieudonné cristalline''. Par ailleurs,

\margin{223}
